\section{Preliminaries and notation}
\label{sec:preliminaries}

\subsection{Polynomial systems, ideals, and varieties}

Let $K$ be a field, let $\overline K$ be its algebraic closure, and write
$R=K[x_1,\ldots,x_n]$. The input is a finite system
$P=\{p_1,\ldots,p_m\}\subseteq R$ of nonzero polynomials. Unless stated
otherwise, a solution is a point in $\overline K^n$. For a finite set
$Q\subseteq R$, its solution set is
\[
    V(Q)=\{a\in\overline K^n:q(a)=0\text{ for every }q\in Q\}.
\]
We call $Q$ satisfiable if $V(Q)\neq\emptyset$ and unsatisfiable (UNSAT) if
$V(Q)=\emptyset$.

The ideal generated by $Q$, denoted $\langle Q\rangle$, consists of all sums
$\sum_{q\in Q}h_q q$ with $h_q\in R$.

For $X\subseteq\overline K^n$, its vanishing ideal $I(X)$ contains every
polynomial in $R$ that vanishes on all of $X$. The radical of an ideal $J$ is
$\sqrt J=\{g\in R:g^r\in J\text{ for some }r\geq1\}$. Hilbert's Nullstellensatz
gives
\[
    I(V(Q))=\sqrt{\langle Q\rangle}.
\]
Thus $V(Q)=\emptyset$ exactly when $\langle Q\rangle=\langle1\rangle$. These
facts connect solution sets with computations in polynomial
ideals~\cite{cox2015ideals,kreuzer2000computational}.

\subsection{Monomial orders and Gr\"obner bases}

Let $\mathbb M_n=\{x^\alpha:\alpha\in\mathbb N^n\}$ be the monomials of $R$.
A monomial order $\prec$ is a total order on $\mathbb M_n$ such that every
nonempty subset has a least element and
\[
    u\prec v \quad\Longrightarrow\quad uw\prec vw
    \qquad (u,v,w\in\mathbb M_n).
\]
For a nonzero polynomial $p$, the leading monomial $\operatorname{LM}(p)$ is
the greatest monomial of $p$ under $\prec$.

For an ideal $J$, define its leading-monomial ideal by
\[
    \operatorname{LM}(J) =\langle\operatorname{LM}(f):f\in J\setminus\{0\}\rangle.
\]
A finite set $G\subseteq J$ is a Gr\"obner basis of $J$ if
\[
    \operatorname{LM}(J) =\langle\operatorname{LM}(g):g\in G\rangle.
\]
$p \in J$ if and only if its division by $G$ has a 0 remainder, so division decides ideal membership.

For $Q=\{q_1,\ldots,q_m\}$, \emph{interreduction} repeatedly applies
\[
    q_i\leftarrow\operatorname{NF}_{Q\setminus\{q_i\}}(q_i),
\]
where $\operatorname{NF}$ denotes the remainder under polynomial division,
until every $q_i$ is unchanged. Since each replacement has the form
$q_i\leftarrow q_i-\sum_{j\ne i}h_jq_j$, it preserves $\langle Q\rangle$ and
therefore $V(Q)$.

Every ideal has a unique reduced Gr\"obner basis for a fixed monomial order.
The reduced basis is $\{1\}$ if and only if $J=\langle1\rangle$. Therefore, a
Gr\"obner basis computation can describe a branch and detect when it is
UNSAT~\cite{cox2015ideals,kreuzer2000computational}.

An UNSAT core of a system $Q$ is a subset $U\subseteq Q$ with $V(U)=\emptyset$.
An UNSAT core is minimal if there is no $U'\subsetneq U$ such that $V(U')=\emptyset$.
Any system containing $U$ is also UNSAT, since adding equations cannot introduce solutions.

\subsection{Factor sets and hitting sets}

Fix a factorization of each input polynomial over $K$:
\[
    p_i=a_i\prod_{f\in C_i}f^{e_{i,f}},
\]
where $a_i\in K\setminus\{0\}$, $C_i$ contains the distinct nonconstant
irreducible factors of $p_i$, and $e_{i,f}\geq1$. Multiplicity does not affect
the zero set. A point satisfies $p_i=0$ exactly when at least one factor in
$C_i$ vanishes at that point.

A hitting set for $C_1,\ldots,C_m$ is a set $H\subseteq\bigcup_i C_i$ such that $H\cap C_i\neq\emptyset$ for every $i$.
It is minimal if no proper subset is also a hitting set.
We write $\mathcal H(P)$ for the family of minimal hitting sets obtained from $P$.
Each $H\in\mathcal H(P)$ defines a branch in which all factors in $H$ vanish~\cite{berge1989hypergraphs}.

\subsection{Encoding nonzero variables}
\label{subsec:nonzero-encoding}

There are two standard encodings of nonzero variables as a polynomial equation.
The first one encodes each $x\neq0$ constraint as
\begin{equation}
\label{eq:single-zeroing-saturation}
    \sigma_x=1-s_x x.
\end{equation}
Here $s_x$ is a new variable.

The second encoding combines the constraints $x\neq0$ for all $x\in N$ into
one equation
\begin{equation}
\label{eq:zeroing-saturation}
    \sigma_N=1-s\prod_{x\in N}x.
\end{equation}

For the first encoding, we project away the auxiliary variables $s_x$. For the
second, we project away $s$. In both cases, the remaining points are the
solutions of the original system that satisfy the nonzero conditions.
